\subsection{Corrected grids without a tree decomposition}\label{sec:balanced}

CT (Algorithm~\ref{alg:ct}) uses its tree decomposition only to compute, on
the current product grid $G=\prod_iG_i$, the corrected minimum
$\beta=\min_GQ$, one minimizer and the min-marginals $m_i(v)$ of $Q=F-D$
(Lemma~\ref{lem:dp}). Any exact method for these quantities can replace the
dynamic program. For the following sign structure, minimum cuts are such a
method, whatever the interaction graph.

\begin{definition}[Balanced quadratic]\label{def:balanced}
A quadratic $F(x)=\frac12x^{\T}Hx+b^{\T}x+c$ is \emph{balanced} if there are
signs $o\in\{-1,1\}^n$ with $o_io_jH_{ij}\le0$ for all $i\ne j$.
\end{definition}

Equivalently, after the substitution $x_i\mapsto o_ix_i$ all off-diagonal
Hessian entries are nonpositive, so $F$ becomes submodular on $\R^n$.
Condition (R2) of Section~\ref{sec:cuts} says that $F$ becomes balanced once
the core coordinates are fixed, and Lemma~\ref{lem:balance} recognizes
balance in linear time.
No sign condition is imposed on the diagonal or on the linear terms. If
$H_{ii}\le0$ for all $i$, then $F$ satisfies~\eqref{eq:cr-class} with an
empty core, and Theorem~\ref{thm:cr-oracle} minimizes it exactly with one
maximum flow.

\begin{lemma}[Grid minimization by one minimum cut]\label{lem:chaincut}
Let $F$ be balanced with signs $o$, let $G_i\subset\Q$ be finite and
nonempty with $\abs{G_i}=k_i+1$, and let $\chi_i:G_i\to\Q$ be arbitrary. Put
$F_\chi(y)=F(y)+\sum_i\chi_i(y_i)$ for $y\in\prod_iG_i$. Then $\min F_\chi$
and a minimizer are obtained from one minimum $\mathsf s$--$\mathsf t$ cut in
a network with $\sum_ik_i+2$ nodes and at most
$3\sum_ik_i+2\sum_{i<j:\,H_{ij}\ne0}k_ik_j$ arcs, whose capacities are
computed from $H$, $b$, the grid nodes and the values of the $\chi_i$ by
polynomially many rational operations. A feasible flow of equal value
certifies the minimum. The same holds for each min-marginal
$\min\{F_\chi(y):y_i=v\}$, with $G_i$ replaced by $\{v\}$.
\end{lemma}

The proof (Appendix~\ref{app:recourse-cuts}) encodes each coordinate by
threshold variables and applies Proposition~\ref{prop:cr-cut}(a); arcs of
infinite capacity exclude the binary vectors that encode no grid node.
Threshold encodings, and the reduction of pairwise problems that are
submodular with respect to label orders to one minimum cut, are classical
\cite[Sections~3--4]{SchlesingerFlach2006}, \cite{Ishikawa2003}.
Lemma~\ref{lem:chaincut} is this construction for quadratic pair terms; what
matters here is that the unary terms $\chi_i$ are arbitrary, so that they
can be the negative corrections $-d_i$ on nonuniform grids.

\begin{theorem}[Balanced box quadratics without a width parameter]\label{thm:balanced}
Let $F$ be a balanced rational quadratic on a mixed box $X$ with input
length $I$. Run CT and EX (Algorithm~\ref{alg:ex}) with $\beta$, a
minimizer and all min-marginals computed by Lemma~\ref{lem:chaincut} instead
of Lemma~\ref{lem:dp}.
\begin{enumerate}[label=(\alph*)]
\item On every instance, CT$(2^{-q})$ halts with a feasible rational point
and a valid path certificate of gap at most $2^{-q}$, in which the dynamic
programs are replaced by minimum cuts with flows of equal value; and EX
halts with a global minimizer and $\OPT$.
\item If $F$ has point growth with condition number $\kappa$, then
CT$(2^{-q})$ uses $(n\kappa(I+q+1))^{O(1)}$ bit operations and EX uses
$(n\kappa(I+1))^{O(1)}$ bit operations.
\end{enumerate}
No decomposition is used, $g$ and $\kappa$ are not inputs, and the validity
of every certificate uses neither growth nor uniqueness.
\end{theorem}

\begin{proof}
(a) In CT and EX, and in the proofs of Propositions~\ref{prop:cellwise}
and~\ref{prop:filter}, Theorem~\ref{thm:certificate},
Lemmas~\ref{lem:inv} and~\ref{lem:states} and
Theorems~\ref{thm:approx}, \ref{thm:transfer} and~\ref{thm:exact}, the tree
decomposition enters only through Lemma~\ref{lem:dp}, which computes
$\beta$, a minimizer of $Q$ and the min-marginals $m_i(v)$ on the current
grid (if $P=\emptyset$, CT minimizes $F$ over the grid
$\prod_i\{\ell_i,u_i\}$, on which $Q=F$). Everything else uses the data of
$F$ directly: the condition of Definition~\ref{def:curvature} holds with
$L_i=\max\{H_{ii},0\}$, because the restriction of $F$ to a coordinate line
has second derivative $H_{ii}$, and REC (Algorithm~\ref{alg:rec}) and
Proposition~\ref{prop:accept} work with $F$ and its gradient. Since $D$ is a
sum of unary terms, $Q=F_\chi$ with $\chi_i=-d_i$, so
Lemma~\ref{lem:chaincut} computes $\beta$ and a minimizer with one cut and
each $m_i(v)$ with one further cut. A flow of equal value certifies each of
these values from below by weak duality, which is what conditions (C1)
and~(C2) of a path certificate require. Hence the validity and termination
statements of these results hold verbatim, including the abort rule of CT
and, under weighted growth, the success of a trial with
$2^\mu\le6\sqrt{\bar\kappa}$; only the operation counts change.

(b) It remains to count cut computations and bit lengths in place of table
operations; the count is in Appendix~\ref{app:recourse-cuts}.
\end{proof}

\begin{corollary}[Balanced after deleting a supplied core]\label{cor:core}
Let $F$ and $X$ be as in Theorem~\ref{thm:balanced}, except that $F$ need not
be balanced. Suppose a set $\mathcal K$ of $k$ coordinates is supplied such
that $o_io_jH_{ij}\le0$ for some signs $o$ and all $i\ne j$ outside
$\mathcal K$. Then Theorem~\ref{thm:balanced} holds for $F$, with the work
bounds of part~(b) multiplied by $K^k$, where $K=O(\sqrt\kappa\log n)$
bounds the nodes per coordinate in every trial that CT runs.
\end{corollary}

\begin{proof}
Fixing the core coordinates at grid nodes changes only unary and constant
terms of the remaining quadratic, which is balanced. Hence $\beta$, a
minimizer and each min-marginal are minima, over at most $K^k$ vectors of
core nodes, of problems solved by Lemma~\ref{lem:chaincut}; a min-marginal
of a core coordinate fixes that coordinate. Every trial that CT runs has at most
$K_\mu\le K_{\mu^*}$ nodes per coordinate, where
$2^{\mu^*}\le6\sqrt{\bar\kappa}\le6\sqrt\kappa$ (proof of
Theorem~\ref{thm:approx}(b) and Appendix~\ref{app:recourse-cuts}), so
$K=K_{\mu^*}=O(\sqrt\kappa\log n)$.
The rest of the proof of Theorem~\ref{thm:balanced} is unchanged.
\end{proof}

Corollary~\ref{cor:core} complements Section~\ref{sec:cuts}: residual
coordinates may now have positive diagonal, but they are gridded and filtered
like the core, so their curvature enters $\kappa$, and point growth of $F$
is needed instead of growth of $V$ in the core coordinates; the residual
graph may still be dense.
Section~\ref{sec:related} compares Theorem~\ref{thm:balanced} with methods
for continuous submodular minimization.
